\documentclass[article]{jss}

\newcommand{\pkgname}{\pkg{AccumulatR}}
\newcommand{\fct}[1]{\code{#1()}}
\newcommand{\class}[1]{\code{#1}}

\author{Niek Stevenson\\Independent researcher}
\Plainauthor{Niek Stevenson}

\title{\pkgname: Structured Evidence-Accumulation Models for Choice and Response-Time Data in \proglang{R}}
\Plaintitle{AccumulatR: Structured Evidence-Accumulation Models for Choice and Response-Time Data in R}
\Shorttitle{\pkgname{} in \proglang{R}}

\Abstract{
  Choice and response-time data are often modeled with evidence-accumulation
  models in which latent finishing times determine both the selected response
  and its latency. Many applications require structured variants of this idea:
  responses may depend on pooled channels, ordered stages, inhibitory processes,
  latent components, shared absence events, or logical combinations of several
  latent completions. \pkgname{} is an \proglang{R}/\proglang{C++} package for
  specifying, simulating, and evaluating likelihoods for structured
  evidence-accumulation models. Its central object is a race architecture rather
  than a single named response-time distribution or estimation routine. Users
  declare accumulators, pools, outcome rules, components, triggers, and mixture
  settings in \proglang{R}; the finalized model object is then used for
  simulation, data preparation, response-probability calculation, and likelihood
  evaluation. Repeated numerical work is delegated to compiled code via
  \pkg{Rcpp}. This paper describes the software design, the mathematical
  representation used by the likelihood evaluator, and reproducible examples
  covering basic races, pooled outcomes, guarded rules, mixtures, and ranked
  responses.
}

\Keywords{evidence accumulation, response times, race models, choice models, \proglang{R}}
\Plainkeywords{evidence accumulation, response times, race models, choice models, R}

\Address{
  Niek Stevenson\\
  Independent researcher\\
  E-mail: \email{niek.stevenson@gmail.com}\\
  URL: \url{https://github.com/niekstevenson/AccumulatR/}
}

\begin{document}

\section{Introduction}
\label{sec:introduction}

Psychologists study decision-making using choice outcomes and the time at which
that choice was made. A model for such data must therefore explain
two things at once: which response occurred and when it occurred.
Evidence-accumulation models (EAMs) do this by introducing latent processes that evolve
over time until a finishing condition is met. The observed response is a label
attached to that finishing event, and the response time is its latency. 

% start changing here
% Diffusion decision models (DDMs), linear ballistic accumulator models, racing diffusion
%models, and lognormal race models differ in their latent dynamics, but share
%this joint modeling strategy \citep{ratcliffDiffusionDecision2016,
% brownSimplestComplete2008, tillmanSequentialSampling2020,
% rouderLognormalRace2015}.

For many data sets, a standard readout is sufficient: one process reaches one
of two boundaries, or the first accumulator to finish determines the response.
Other experiments make the readout part of the hypothesis. A response may be
produced by the first member of a pool, by the completion of several processes,
by a target process that must finish before an inhibiting process, by a later
stage that can start only after an earlier one, or by a latent processing
component selected on that trial. These are substantive claims about cognitive
architecture, but they are also statistical claims. Changing the readout changes
the event whose probability is evaluated and therefore changes the likelihood of
the observed choice and response time.

The current software ecosystem is organized mainly around named distributions
and fitting workflows. Packages such as \pkg{RWiener}, \pkg{fddm},
\pkg{ddModel}, and \pkg{ream} provide distributional calculations for DDM or
two-threshold models \citep{wabersichRWienerPackage2014,
fosterFddmPackage2024, linDdmodelPackage2025, hartmannReamPackage2025}.
Packages such as \pkg{rtdists}, \pkg{lbaModel}, \pkg{CircularDDM},
\pkg{dynConfiR}, \pkg{eam}, \pkg{ssms}, and
\pkg{SequentialSamplingModels.jl} provide density, simulation, likelihood, or
prediction functions for accumulator, race, continuous-response, confidence, or
broader sequential-sampling model families \citep{singmannRtdistsPackage2026,
linLbaModelPackage2025, linCircularDDMPackage2018,
hellmannDynconfirPackage2025, zhuEamPackage2026,
fenglerLikelihoodApproximation2021, fernandezSequentialSamplingModels2024}.
Estimation systems such as \pkg{fast-dm}, \pkg{PyDDM}, \pkg{PyBEAM},
\pkg{dRiftDM}, \pkg{DMCfun}, \pkg{diffIRT}, \pkg{RegDDM},
\pkg{DMC}/\pkg{ggdmc}, \pkg{EMC2}, \pkg{pmwg}, \pkg{brms}, \pkg{bmm},
\pkg{HDDM}, and \pkg{HSSM} provide optimization, Bayesian, hierarchical,
regression, simulator-based, or amortized-likelihood workflows
\citep{vossFastdmFree2007, shinnPyDDM2020, murrowPyBEAM2023,
koobDriftdmPackage2026, mackenzieDmcfunPackage2024, molenaarDiffIRT2015,
jinRegddmPackage2025, heathcoteDynamicModels2019, linGgdmcPackage2025,
stevensonEMC2Package2024, cooperPmwgPackage2024, burknerBrmsPackage2017,
popovBmmPackage2026, wieckiHDDM2013, fenglerHssmPackage2026}.
The resulting ecosystem is valuable, but its reusable objects are usually
distributions, simulators, fitted models, or sampling workflows. The readout
architecture is typically fixed by the named model family. When that
architecture is the object to be varied, the user often has to force the
hypothesis into an existing readout or write a custom simulator and likelihood.

This distinction has a long history in reaction-time theory. Serial--parallel
systems, critical-path models, and queueing-network accounts all treat process
architecture as part of the model rather than as verbal interpretation after
fitting \citep{schweickertCriticalPath1978, townsendSerialparallelDilemma2004,
liuQueueingNetwork2007}. Race models are a useful way to bring that idea into
evidence accumulation because they represent a trial by several latent finishing
times \citep{heathcoteWinnerTakes2022}. The usual winner-takes-all race is then
only the simplest order constraint. Pooled, staged, inhibited, delayed, mixed,
or ranked responses can be represented by different constraints on the same
kind of latent objects.


\pkgname{} addresses this gap by making the race architecture a reusable
software object. A model specification contains the latent accumulators and the
rules that map their finishing times to observed responses: pools, logical
outcome rules, triggers, components, mixtures, onsets, and ranked readouts. The
same finalized object is used for simulation, response-probability calculation,
data preparation, and likelihood evaluation. Internally, the evaluator lowers
outcome rules to order constraints over latent finishing times and reuses
compiled model state across parameter evaluations. The package therefore
provides a model-evaluation layer for structured evidence-accumulation models;
it does not try to replace optimizers, samplers, or hierarchical modeling
systems.

The examples in this paper are drawn from cognitive tasks because such tasks
motivate structured response architectures. The software contribution, however,
is statistical: \pkgname{} gives users a compact \proglang{R} interface for
declaring those architectures and an \proglang{R}/\proglang{C++} evaluator for
the probabilities they imply \citep{R, Eddelbuettel+Francois:2011}.
Section~\ref{sec:design} introduces the package objects and public API.
Section~\ref{sec:model} summarizes the mathematical representation.
Section~\ref{sec:workflow} gives the basic \proglang{R} workflow.
Section~\ref{sec:examples} presents structured model examples.
Sections~\ref{sec:likelihood} and~\ref{sec:implementation} describe likelihood
evaluation and implementation details, followed by a comparison to related
software and a discussion.

\section{Software design}
\label{sec:design}

\subsection{Objects and workflow}

The package distinguishes model structure from numerical parameter values. A
\class{race\_spec} object describes the architecture: accumulators, pools,
observed outcome rules, components, triggers, mixture settings, and parameter
naming rules. A finalized \class{model\_structure} object is produced for
simulation, data preparation, response probabilities, and likelihood contexts.
The compiled likelihood context is built once and then reused across many
parameter evaluations.

\begin{table}[ht!]
\centering
\begin{tabular}{@{}p{2.0cm}p{5.2cm}p{6.1cm}@{}}
\hline
Stage & Main functions & Role \\ \hline
Spec. &
\fct{race\_spec}, \fct{add\_accumulator},\newline \fct{add\_pool},
\fct{add\_outcome} &
Declare the latent accumulators, pooled sources, and observed response labels. \\
Rules &
\fct{first\_of}, \fct{all\_of}, \fct{none\_of},\newline
\fct{inhibit}, \fct{after}, \fct{add\_trigger} &
Express logical outcome rules, blocking rules, staged onsets, and shared
absence events. \\
Params &
\fct{set\_parameters}, \fct{par\_names},\newline
\fct{build\_param\_matrix} &
Control public parameter names and expand named values to trial-wise parameter
rows. \\
Mixing &
\fct{add\_component}, \fct{set\_mixture} &
Define component membership and specify whether mixture weights are fixed or
estimated as sampled-mixture parameters. \\
Runtime &
\fct{finalize\_model}, \fct{simulate},\newline
\fct{prepare\_data}, \fct{make\_context},\newline
\fct{log\_likelihood} &
Finalize a specification, simulate data, prepare observations, and evaluate the
likelihood. \\
Inspect &
\fct{processing\_tree},\newline \fct{response\_probabilities},
\fct{complexity\_metrics} &
Inspect response logic, marginal response probabilities, and compiled
complexity diagnostics. \\ \hline
\end{tabular}
\caption{\label{tab:api} Main public functions in \pkgname{}.}
\end{table}

The workflow in Table~\ref{tab:api} is deliberately explicit. Model-building
functions return modified specifications, so they can be chained with the base
\proglang{R} pipe. Parameter values are not stored in the model object. Instead
\fct{build\_param\_matrix} expands a named vector into the rectangular or long
parameter layout expected by the simulator and likelihood evaluator. This
separation lets the same model architecture be evaluated under many trial-wise
parameter sets.

\subsection{Parameter naming}

Parameter names are public interface, not an implementation detail. By default,
compatible accumulator parameters are grouped. For example, a two-accumulator
lognormal race has public parameters \code{m}, \code{s}, and \code{t0}; both
accumulators share those values unless the specification requests otherwise.
The function \fct{set\_parameters} changes this mapping. Its
\code{separate} argument splits grouped parameters, \code{share} intentionally
maps several internal parameters to one new public name, and \code{rename}
changes a public name without merging distinct parameters. The result can be
inspected with \fct{par\_names}.

This design avoids the common ambiguity between an accumulator-level parameter
and a public parameter. The name \code{left.m} is valid only if the
specification has separated \code{m} for the accumulator \code{left}. Without
that declaration, the valid name is \code{m}. The same mechanism is used for
models that combine distribution families. If the same suffix has incompatible
meanings across distributions, the default public name is qualified by
distribution, for example \code{lognormal.s} and \code{RDM.s}.

\subsection{Components and mixtures}

Components represent qualitatively different process variants within one model.
A component declaration defines membership only: it says which accumulators are
active in a component, and optionally how many ordered responses that component
records. Mixture probabilities are configured separately with
\fct{set\_mixture}. In \code{mode = "fixed"}, component weights are supplied as
known probabilities. In \code{mode = "sample"}, non-reference components expose
probability parameters such as \code{p.fast}; the reference receives the
residual probability.

The data determine whether a component is observed on a trial. If a
\code{component} column is present, likelihood evaluation conditions on that
component. If the column is absent or contains missing component labels, the
likelihood marginalizes over components using the model's mixture
specification. The same convention is used by \fct{response\_probabilities}.

\section{Model representation}
\label{sec:model}

\subsection{Leaves, sources, and response rules}

The basic latent object is a finishing time. For a leaf accumulator $a$, let
$T_a$ denote its finishing time, with distribution function $F_a$, density
$f_a$, and survival function $S_a$:
%
\begin{equation}
  F_a(t) = \Pr(T_a \leq t), \qquad
  f_a(t) = F_a'(t), \qquad
  S_a(t) = 1 - F_a(t).
  \label{eq:leaf}
\end{equation}
%
A source is either a leaf or a pool. An outcome rule defines when an observed
response becomes available. The simplest outcome rule is a direct readout from
one source. More complex rules use first-of, all-of, absence, and inhibition
operators.

\begin{table}[t!]
\centering
\begin{tabular}{lll}
\hline
Distribution & Public parameters & Notes \\ \hline
\code{lognormal} & \code{m}, \code{s} & Log-scale location and spread. \\
\code{gamma} & \code{shape}, \code{rate} & Gamma finishing-time distribution. \\
\code{exgauss} & \code{mu}, \code{sigma}, \code{tau} & Gaussian plus exponential finishing time. \\
\code{LBA} & \code{v}, \code{B}, \code{A}, \code{sv} & Linear ballistic accumulator. \\
\code{RDM} & \code{v}, \code{B}, \code{A}, \code{s} & Race diffusion model parameterization. \\ \hline
\end{tabular}
\caption{\label{tab:distributions} Accumulator distributions currently
available in \pkgname{}. All accumulators also have the non-decision-time
parameter \code{t0}, which defaults to zero if omitted.}
\end{table}

For a simple independent race in which response $r$ is observed at time $t$,
the familiar density is
%
\begin{equation}
  \ell_r(t) = f_r(t) \prod_{j \neq r} S_j(t).
  \label{eq:simple-race}
\end{equation}
%
\pkgname{} treats Equation~\ref{eq:simple-race} as a special case of a more
general order-region representation. A target outcome may become true through
several transition scenarios, and each scenario is combined with the condition
that no competing observed outcome has occurred earlier. For direct,
independent races this lowers immediately to products of densities and
survival functions. For shared, nested, or guarded rules the evaluator keeps
the joint order constraints until they can be projected correctly.

\subsection{Pools and logical rules}

A pool combines several sources into one source. If pool $P$ has members
$T_1,\ldots,T_n$ and finishes when the $k$th member finishes, then
%
\begin{equation}
  F_P(t) =
  \Pr\left(\sum_{j=1}^{n} 1\{T_j \leq t\} \geq k\right).
  \label{eq:pool}
\end{equation}
%
For independent exchangeable members with common $F$ and $S = 1 - F$, this
reduces to
%
\begin{equation}
  F_P(t) = \sum_{m=k}^{n} {n \choose m} F(t)^m S(t)^{n-m}.
  \label{eq:pool-iid}
\end{equation}
%
The software does not require exchangeability; the pool is compiled as order
constraints over its members.

The helpers \fct{first\_of} and \fct{all\_of} define disjunction and
conjunction over response rules. With independent children, \fct{first\_of}
becomes true when the earliest child becomes true, and \fct{all\_of} becomes
true when the last required child becomes true:
%
\begin{equation}
  f_{\mathrm{first}}(t)
  =
  \sum_{r=1}^{m} f_r(t) \prod_{q \neq r} S_q(t),
  \qquad
  f_{\mathrm{all}}(t)
  =
  \sum_{r=1}^{m} f_r(t) \prod_{q \neq r} F_q(t).
  \label{eq:first-all}
\end{equation}
%
The helper \fct{none\_of} is an absence condition. The helper \fct{inhibit}
defines a response-generating rule and a blocker. For independent direct
sources,
%
\begin{equation}
  f_{\mathrm{inhibit}(A,B)}(t) = f_A(t) S_B(t),
  \label{eq:inhibit}
\end{equation}
%
but the implemented representation also handles composite blockers and shared
sources.

\subsection{Onsets, triggers, mixtures, and ranks}

An accumulator can have an absolute onset or a chained onset. With an absolute
delay $d_a$, the observed finishing time is $T_a = d_a + T_a^\star$. With
\fct{after}, the accumulator begins only after another source has finished. If
accumulator $a$ starts after source $U$ with lag $\delta$, then
%
\begin{equation}
  T_a = T_U + \delta + T_a^\star.
  \label{eq:after}
\end{equation}
%
The delayed accumulator then participates in the same outcome rules as any
other accumulator.

Triggers control whether accumulators are present on a trial. A trigger is a
named absence-probability parameter. All members in one \fct{add\_trigger}
call share the same absence draw; separate calls create independent absence
draws. If a shared trigger has absence probability $q_g$, the trial likelihood
is a mixture over the trigger-on and trigger-off states:
%
\begin{equation}
  L = (1 - q_g) L^{\mathrm{on}} + q_g L^{\mathrm{off}}.
  \label{eq:trigger}
\end{equation}
%
Mixture components use the same likelihood principle. If a component is
observed, the trial likelihood conditions on that component. If it is latent,
the likelihood averages over components:
%
\begin{equation}
  L_i = \sum_c \pi_{ic} L_{ic}.
  \label{eq:mixture}
\end{equation}
%
For ranked-response data, the model can retain more than one ordered response
per trial by setting \code{n\_outcomes} in \fct{race\_spec}. The likelihood is
then evaluated sequentially over the observed response-time pairs
$(r_1,t_1),\ldots,(r_K,t_K)$.

\section[Using AccumulatR]{Using \pkgname{}}
\label{sec:workflow}

\subsection{A two-choice race}

The following example defines a two-choice lognormal race. Before parameter
separation, both accumulators share the default public parameters:

\begin{CodeChunk}
\begin{CodeInput}
R> library("AccumulatR")
R> spec <- race_spec() |>
+   add_accumulator("left", "lognormal") |>
+   add_accumulator("right", "lognormal") |>
+   add_outcome("left", "left") |>
+   add_outcome("right", "right")
R> par_names(spec)
\end{CodeInput}
\begin{CodeOutput}
[1] "m"  "s"  "t0"
\end{CodeOutput}
\end{CodeChunk}

To give the two accumulators separate lognormal locations and spreads, the
specification uses \fct{set\_parameters}. The non-decision time \code{t0}
remains shared and is not supplied below, so it defaults to zero.

\begin{CodeChunk}
\begin{CodeInput}
R> model <- spec |>
+   set_parameters(separate = list(m = TRUE, s = TRUE)) |>
+   finalize_model()
R> pars <- c(
+   left.m = log(0.28), left.s = 0.16,
+   right.m = log(0.35), right.s = 0.18
+ )
R> param_df <- build_param_matrix(model, pars, n_trials = 120)
R> sim <- simulate(model, param_df, seed = 20260612)
R> head(sim, 4)
\end{CodeInput}
\begin{CodeOutput}
  trials     R        rt
1      1 right 0.3167664
2      2  left 0.2847782
3      3  left 0.3259237
4      4  left 0.2515444
\end{CodeOutput}
\end{CodeChunk}

The simulated data use the same columns expected for real behavioral data:
\code{trials}, \code{R}, and \code{rt}. The likelihood workflow prepares the
data once, builds a compiled context once, and evaluates parameter matrices as
needed.

\begin{CodeChunk}
\begin{CodeInput}
R> prepared <- prepare_data(model, sim[c("trials", "R", "rt")])
R> ctx <- make_context(model)
R> log_likelihood(ctx, prepared,
+   build_param_matrix(model, pars, trial_df = prepared))
\end{CodeInput}
\begin{CodeOutput}
[1] 157.0651
\end{CodeOutput}
\end{CodeChunk}

Response probabilities can be computed without simulation:

\begin{CodeChunk}
\begin{CodeInput}
R> one_trial <- build_param_matrix(model, pars, n_trials = 1)
R> round(response_probabilities(model, one_trial), 3)
\end{CodeInput}
\begin{CodeOutput}
 left right
0.823 0.177
\end{CodeOutput}
\end{CodeChunk}

The value returned by \fct{log\_likelihood} is a density-based log-likelihood,
so it need not be negative. With continuous response-time models, densities can
be greater than one.

\subsection{Trial-wise parameters}

The parameter matrix is trial-wise by construction. A typical fitting function
therefore takes a low-dimensional vector of free parameters, maps it to the
public parameter names of the model, and then calls \fct{build\_param\_matrix}
with the prepared data. For example, a minimal negative log-likelihood
function for the two-choice race is:

\begin{Code}
neg_loglik <- function(theta) {
  est <- c(
    left.m = theta[["left.m"]],
    left.s = exp(theta[["log_left.s"]]),
    right.m = theta[["right.m"]],
    right.s = exp(theta[["log_right.s"]])
  )
  params <- build_param_matrix(model, est, trial_df = prepared)
  -as.numeric(log_likelihood(ctx, prepared, params))
}
\end{Code}

The package does not impose an optimizer or sampler. The likelihood can be
used with \fct{optim}, custom maximum-likelihood routines, or external
Bayesian engines that call back into \proglang{R}.

\section{Structured model examples}
\label{sec:examples}

\subsection{Pooled evidence}

Pools represent responses generated by a threshold number of latent channels.
In the following specification, response \code{A} is observed when any two of
three \code{A} channels have finished, while response \code{B} is generated by
a single competing accumulator.

\begin{Code}
pool_model <- race_spec() |>
  add_accumulator("A1", "lognormal") |>
  add_accumulator("A2", "lognormal") |>
  add_accumulator("A3", "lognormal") |>
  add_accumulator("B", "lognormal") |>
  add_pool("A_pool", c("A1", "A2", "A3"), k = 2L) |>
  add_outcome("A", "A_pool") |>
  add_outcome("B", "B") |>
  set_parameters(separate = list(m = TRUE, s = TRUE)) |>
  finalize_model()
\end{Code}

When \code{A} is observed, its response time is the finishing time of the
second finisher in \code{A_pool}. The likelihood uses that pooled completion
rule directly; the user does not need to derive a separate density for the
order statistic.

\subsection{Logical response rules}

Observed responses can be defined by logical rules rather than by single
accumulators. The next model says that a response occurs only if an encoding
stage has finished, either choice route has finished, and the stop process has
not occurred.

\begin{Code}
logical_model <- race_spec() |>
  add_accumulator("encoding", "lognormal") |>
  add_accumulator("left_route", "lognormal") |>
  add_accumulator("right_route", "lognormal") |>
  add_accumulator("stop", "lognormal") |>
  add_outcome(
    "choice",
    all_of(
      "encoding",
      first_of("left_route", "right_route"),
      none_of("stop")
    )
  ) |>
  finalize_model()
\end{Code}

For stop-like architectures where one process generates a response and another
can block it, \fct{inhibit} gives a more direct expression:

\begin{Code}
stop_model <- race_spec() |>
  add_accumulator("go", "lognormal") |>
  add_accumulator("stop", "lognormal") |>
  add_outcome("respond", inhibit("go", "stop")) |>
  finalize_model()
\end{Code}

Both examples use string labels directly. The helper
\fct{build\_outcome\_expr} is available for programmatic builders that need
quoted \proglang{R} expressions.

\subsection{Chained onsets and shared triggers}

Chained onsets express staged processing. Here accumulator \code{C} starts
only after latent accumulator \code{B} has finished:

\begin{Code}
onset_model <- race_spec() |>
  add_accumulator("A", "lognormal") |>
  add_accumulator("B", "lognormal") |>
  add_accumulator("C", "lognormal", onset = after("B")) |>
  add_outcome("A", "A") |>
  add_outcome("C", "C") |>
  set_parameters(separate = list(m = TRUE, s = TRUE)) |>
  finalize_model()
\end{Code}

Triggers express shared absence events. In the next model, the two \code{go}
accumulators share one absence draw and one absence-probability parameter:

\begin{Code}
trigger_model <- race_spec() |>
  add_accumulator("go1", "lognormal") |>
  add_accumulator("go2", "lognormal") |>
  add_outcome("R1", "go1") |>
  add_outcome("R2", "go2") |>
  add_trigger("shared_trigger", members = c("go1", "go2")) |>
  set_parameters(separate = list(m = TRUE, s = TRUE)) |>
  finalize_model()
\end{Code}

If separate \fct{add\_trigger} calls are used instead, the absence draws are
independent. Their probabilities can still be given a shared public parameter
through \fct{set\_parameters}; that changes parameter naming but not the
independence of the trigger draws.

\subsection{Latent mixtures}

Components are useful when trials may come from qualitatively different
processing modes. The following model has a fast and a slow target route, with
one competitor. The component declarations define which accumulators are active
in each component, and \fct{set\_mixture} declares the component weights to be
sampled parameters.

\begin{Code}
mix_model <- race_spec() |>
  add_accumulator("target_fast", "lognormal") |>
  add_accumulator("target_slow", "lognormal") |>
  add_accumulator("competitor", "lognormal") |>
  add_pool("target", c("target_fast", "target_slow")) |>
  add_outcome("R1", "target") |>
  add_outcome("R2", "competitor") |>
  add_component("fast", members = c("target_fast", "competitor")) |>
  add_component("slow", members = c("target_slow", "competitor")) |>
  set_parameters(separate = list(m = TRUE, s = TRUE)) |>
  set_mixture(mode = "sample", reference = "slow") |>
  finalize_model()
\end{Code}

For this model, the public parameter \code{p.fast} is the probability of the
fast component and the slow component receives the residual mass. If the data
include a \code{component} column, the likelihood conditions on the observed
component. Without that column, \code{fast} and \code{slow} are treated as
latent processing modes and marginalized.

\subsection{Ranked responses}

Some experiments record more than one ordered response on a trial. Setting
\code{n\_outcomes = 2L} asks the simulator and likelihood evaluator to retain
the first and second finishing responses:

\begin{Code}
rank_model <- race_spec(n_outcomes = 2L) |>
  add_accumulator("A", "lognormal") |>
  add_accumulator("B", "lognormal") |>
  add_outcome("A", "A") |>
  add_outcome("B", "B") |>
  set_parameters(separate = list(m = TRUE, s = TRUE)) |>
  finalize_model()
\end{Code}

The corresponding data columns are \code{R}, \code{rt}, \code{R2}, and
\code{rt2}. Higher ranks follow the same naming pattern.

\section{Likelihood evaluation and fitting}
\label{sec:likelihood}

The likelihood interface is intentionally narrow. \fct{prepare\_data} checks
and reshapes behavioral observations. \fct{make\_context} compiles the
model-side runtime state. The likelihood call combines a context, prepared
data, and a parameter matrix. This separation matters for estimation because
model compilation is usually more expensive than evaluating one new parameter
matrix.

\begin{Code}
prepared <- prepare_data(model, data_df)
ctx <- make_context(model)

objective <- function(theta) {
  values <- transform_theta_to_named_parameters(theta)
  params <- build_param_matrix(model, values, trial_df = prepared)
  -as.numeric(log_likelihood(ctx, prepared, params))
}

fit <- optim(start, objective, method = "Nelder-Mead")
\end{Code}

The transformation placeholder in this example is user code: it maps
unconstrained optimization variables to the public parameter names of the
model. For example, positive spread, boundary, and non-decision-time parameters
are commonly optimized on the log scale and exponentiated before calling
\fct{build\_param\_matrix}. This keeps the package focused on model evaluation
while leaving estimation strategy to the analysis.

The parameter matrix may be rectangular, with one row per accumulator per
trial, or long, with rows aligned to prepared data for component-specific
architectures. The \code{layout = "auto"} default selects the appropriate
layout for most uses. Prepared data can also be compressed when repeated
observations share the same prepared representation; in that case the
likelihood sums weighted compact rows and expands trial-wise values when
requested.

\section{Implementation}
\label{sec:implementation}

The implementation has three layers. The first layer is the \proglang{R}
specification interface. It validates names, stores accumulator definitions,
records outcome expressions, normalizes parameter mappings, and finalizes a
model into a runtime-oriented structure. The second layer prepares data and
parameters for repeated evaluation. The third layer is the compiled evaluator,
which receives the finalized model plan and performs simulation or likelihood
calculations.

The exact likelihood evaluator lowers outcome rules into transition scenarios.
A scenario records which latent source releases a target rule at the observed
time, which other sources must already be available, and which blockers or
competitors must not have occurred earlier. These statements are represented as
order constraints over latent finishing times. The compiler simplifies
contradictory or redundant cells, projects simple constraints to density,
distribution, or survival factors, and retains numerical integration only where
the order structure cannot be reduced further.

This representation is important for models with shared sources. If two
observed outcomes both depend on a common latent stage, they can tie with
positive density because the same latent finishing time releases both rules.
Treating every tie as a zero-probability boundary would be incorrect. By
keeping shared latent identity in the region representation, \pkgname{} can
distinguish ordinary zero-measure equality from model-induced equality.

Most computational cost that depends only on model structure is paid when the
context is built. Repeated likelihood calls reuse the compiled plan and mostly
pay the cost of evaluating densities, probabilities, and remaining integrals
under new parameter values. For complicated architectures the number of
symbolic cells can grow quickly. When a context is built with
\code{diagnostics = TRUE}, \fct{complexity\_metrics} reports the compiled
complexity so the user can inspect whether a specification is becoming
computationally expensive.

\section{Relation to existing software}
\label{sec:comparison}

\pkgname{} is not a replacement for specialized diffusion-model packages,
model-specific fitters, or full hierarchical estimation systems. The relevant
software landscape is broad, and several packages span more than one role. For
the present comparison the useful distinction is between the statistical object
handled by a package and the level of workflow it exposes.

The distinguishing feature of \pkgname{} is not that it is more general than
all of these systems. It is more specific. Standard race models already provide
a flexible alternative to two-threshold diffusion models because they represent
several latent finishing times rather than a single relative-evidence process.
However, most software for race models still treats the mapping from latent
finishing times to observed responses as part of the named model: for example,
the first accumulator to finish wins, or a fixed LBA/RDM/lognormal-race
readout is assumed. \pkgname{} makes that readout architecture explicit. The
same object can represent ordinary races, pooled accumulators, staged routes,
latent mixtures, guarded responses, shared absence events, and ranked
responses, and the same object is used for simulation, inspection,
response-probability calculation, and likelihood evaluation.

There are corresponding limitations. The package currently assumes that leaf
accumulators are conditionally independent given parameters, components, and
trigger states. It supplies likelihood evaluation but not a complete
hierarchical modeling language. Users who need group-level priors, posterior
sampling diagnostics, or formula-based regression over parameters must supply
that estimation layer separately. In exchange, they obtain a model object whose
architecture can be manipulated directly in \proglang{R}.

Table~\ref{tab:software-comparison} summarizes this classification. The rows
classify packages by their most relevant role for the present comparison, not
by every feature they expose. Some packages span categories; for example,
\pkg{HSSM} supports DDMs and broader SSMs, but is listed with broad EAM
estimation because its current user-facing role is hierarchical estimation with
simulator and amortized-likelihood infrastructure.

\begin{table}[ht!]
\centering
\small
\setlength{\tabcolsep}{3pt}
\begin{tabular}{@{}p{2.8cm}p{4.7cm}p{6.1cm}@{}}
\hline
Category & Examples & Relation to \pkgname{} \\ \hline
DDM/two-threshold evaluation &
\pkg{RWiener}, \pkg{fddm}, \pkg{ddModel}, \pkg{ream}, \pkg{PyDDM} &
Evaluate or simulate named DDM, generalized DDM, or two-threshold accumulation
processes. The response mapping is usually fixed by the two-boundary process. \\
Race/broader evaluation &
\pkg{rtdists}, \pkg{lbaModel}, \pkg{CircularDDM}, \pkg{dynConfiR},
\pkg{eam}, \pkg{ssms}, \pkg{SequentialSamplingModels.jl} &
Provide densities, simulators, or prediction functions for LBA, race,
continuous-response, confidence, or broader SSM families. The response rule is
usually the one built into the named model. \\
DDM/two-threshold estimation &
\pkg{fast-dm}, \pkg{dRiftDM}, \pkg{DMCfun}, \pkg{diffIRT}, \pkg{RegDDM},
\pkg{brms}, \pkg{bmm}, \pkg{HDDM}, \pkg{PyBEAM}, \pkg{PyDDM} &
Estimate DDMs, two-threshold models, task-specific DDM extensions, or
hierarchical regressions. \\
Race/broad EAM estimation &
\pkg{glba}, \pkg{DMC}/\pkg{ggdmc}, \pkg{EMC2}, \pkg{pmwg}, \pkg{eam},
\pkg{HSSM} &
Estimate LBA, RDM, lognormal-race, and other SSM families, often using
hierarchical, simulator-based, or amortized-likelihood workflows. \pkgname{}
supplies the architecture-specific likelihood object that such workflows
otherwise need to hand-code. \\
\hline
\end{tabular}
\caption{\label{tab:software-comparison} Related software for evidence
accumulation and sequential-sampling models, grouped by model family and
workflow role.}
\end{table}

\section{Replication material}
\label{sec:replication}

The manuscript directory contains executable \proglang{R} scripts under
\code{dev/LatexFiles/examples}. The script \code{01-basic-workflow.R}
reproduces the two-choice race in Section~\ref{sec:workflow}, including the
reported response probabilities. The script \code{02-structured-models.R}
constructs pooled, logical, trigger, onset, and ranked-response examples. The
script \code{03-mixtures-and-fitting.R} constructs the sampled-mixture example
and runs a small likelihood-based estimate of the mixture weight. The top-level
\code{code.R} script sources the examples and prints \fct{sessionInfo}, which
is the form expected for JSS replication material.

\section{Discussion}
\label{sec:discussion}

\pkgname{} provides a declarative interface for a class of structured
evidence-accumulation models. Its core design choice is to make architecture
explicit: accumulators, pools, outcome rules, triggers, components, and ranked
readouts are part of the model object rather than ad hoc code around a fixed
density. The same object then supports simulation, response-probability
calculation, data preparation, and likelihood evaluation.

This architecture-first design is most useful for analyses where the main
scientific hypothesis concerns process structure. It also keeps the package
honest about what it does not do. \pkgname{} evaluates models; it does not
attempt to be a universal optimizer, sampler, or hierarchical modeling system.
That boundary keeps the interface small and makes it easier to combine the
package with external estimation tools.

The most important future work is broader empirical demonstration and a richer
estimation interface around the existing likelihood engine. Additional
finishing-time families, more diagnostic summaries, and tighter integration
with hierarchical workflows would also be natural extensions. The present
software already supplies the central machinery needed for those extensions: a
compact model language for structured accumulator architectures and a reusable
compiled evaluator for the likelihoods implied by those architectures.

\section*{Acknowledgments}

The package and manuscript benefited from the published literature on
sequential-sampling models and from the \emph{Journal of Statistical Software}
style and replication guidelines.

\bibliography{../refs,accumulatr-jss-draft}

\end{document}
